\begin{landscape}
    \begin{table}
        \centering
        \begin{talltblr}[
            caption={Komparasi Penelitian Terdahulu yang Berkenaan dengan Topik LaTeX (Hanya Contoh)}, 
            entry={Komparasi Penelitian Terdahulu yang Berkenaan dengan Topik LaTeX (Hanya Contoh)},
            label={table:komparasi-penelitian}
            ]{
                colspec={c X X X X X X}, 
                cells={font=\footnotesize},
                hspan=minimal, 
                hline{1,2,Z}={solid}, 
                row{1}={font=\footnotesize\bfseries}
            }
            No. & Peneliti \& Tahun & Judul Penelitian & Metode \& Sasaran & Persamaan & Perbedaan & Temuan Utama \\
            
            % Baris 1
            1 & {Sapitri \textit{et al}. \\ 2024} &
            Pelatihan Penggunaan LaTeX untuk Penulisan Skripsi bagi
            Mahasiswa Program Studi Sarjana Matematika FMIPA
            Universitas Lampung &
            Pelatihan interaktif \& praktik langsung berbasis platform Overleaf. Sasaran: Mahasiswa S1 Matematika FMIPA Unila. &
            Membahas materi pelatihan LaTeX via Overleaf untuk dokumen ilmiah dan menyediakan \textit{template} format resmi. &
            Sasaran pelatihan adalah mahasiswa S1, fokusnya pada penulisan skripsi akademik mahasiswa. &
            Mahasiswa memperoleh pengetahuan teknis langsung dan \textit{template} skripsi yang sesuai pedoman akademik. \\
            
            % Baris 2
            2 & {Fitriani \textit{et al.} \\ 2024} & 
            Pelatihan LaTeX Menggunakan Overleaf untuk Meningkatkan Kemampuan Penulisan Karya Ilmiah bagi Dosen di Pringsewu &
            Pelatihan, pendampingan, dan evaluasi pengisian kuesioner. Sasaran: Para dosen di Pringsewu. &
            Sama-sama melatih software pengolah dokumen LaTeX (Overleaf) untuk meningkatkan kualitas karya ilmiah. &
            Sasaran pelatihan adalah dosen, lokasinya di Pringsewu, fokusnya peningkatan publikasi ilmiah dosen. &
            Meningkatkan pemahaman teoritis dan kemampuan praktis dosen dalam manajemen tata letak jurnal/buku.
        \end{talltblr}
    \end{table}
\end{landscape}